\section{Conception de la maquette}\label{sec:maquette}

La maquette est orchestrée avec Docker Compose.
Elle comporte les services interconnectés exigés par le sujet : l'application web DVWA et son service de données MariaDB\@.
Deux composants supplémentaires séparent les rôles d'attaque, de détection et de filtrage : un conteneur Kali et une passerelle.

\begin{figure}[H]
    \centering
    \begin{tikzpicture}[
        node distance=1.1cm,
        box/.style={draw, rounded corners, align=center, minimum width=2.6cm, minimum height=1.05cm, fill=black!4},
        flow/.style={-{Stealth[length=2.2mm]}, thick}
    ]
        \node[box] (kali) {Kali Linux\\172.28.10.10\\\scriptsize réseau \texttt{attack}};
        \node[box, right=of kali] (gateway) {Passerelle\\Suricata + iptables\\\scriptsize 172.28.10.254 / 172.28.20.254};
        \node[box, right=of gateway] (dvwa) {DVWA\\172.28.20.20\\\scriptsize réseau \texttt{app}};
        \node[box, below=of dvwa] (db) {MariaDB\\172.28.30.30\\\scriptsize réseau \texttt{data}};
        \draw[flow] (kali) -- node[above, font=\scriptsize]{HTTP} (gateway);
        \draw[flow] (gateway) -- (dvwa);
        \draw[flow] (dvwa) -- (db);
    \end{tikzpicture}
    \caption{Architecture logique de la maquette isolée.}\label{fig:architecture}
\end{figure}

DVWA publie son interface visuelle seulement sur l'interface de boucle locale de l'hôte (port 8080).
Cet accès facilite l'observation de l'application, mais contourne volontairement la passerelle et n'est jamais utilisé pour les scénarios d'attaque.
Le poste Kali ne possède qu'un accès routé vers le réseau applicatif, obligatoirement au travers de la passerelle.
MariaDB n'est raccordée qu'au réseau \texttt{data} : elle n'est donc pas joignable directement depuis Kali.
Ce cloisonnement limite le périmètre de l'exercice tout en conservant des flux représentatifs d'une architecture web à plusieurs tiers.

\begin{table}[H]
    \centering
    \caption{Rôle des composants de la maquette.}\label{tab:composants}
    \begin{tabular}{@{}p{0.22\textwidth}p{0.27\textwidth}p{0.39\textwidth}@{}}
        \toprule
        Composant & Fonction & Justification \\
        \midrule
        Kali & Poste d'audit et de test & Centralise les outils employés dans le périmètre autorisé. \\
        Passerelle & Routage, IDS et filtrage & Observe les requêtes avant la traduction d'adresse et contrôle les flux. \\
        DVWA & Application cible & Fournit un environnement explicitement vulnérable et reproductible~\cite{dvwa}. \\
        MariaDB & Service de données & Représente le composant contenant les données de l'application. \\
        \bottomrule
    \end{tabular}
\end{table}
